\chapter{Glossary}
\label{app:glossary}

The terms the book uses, from language models and agents, from biodiversity
data and from the statistics of testing a model. Protein and genomic language
models, and the vocabulary of sequence-based biodiversity work, are outside
the book's scope and are not defined here.

\begin{description}[itemsep=0.2em]

\item[Access and benefit-sharing (ABS)] The terms under the Convention on
  Biological Diversity on which a user of genetic resources obtains the
  providing country's consent and shares the benefits
  \parencite{scholz2022dsi} (\Cref{sec:abs}).

\item[Active parameters] See \emph{Mixture-of-experts}.

\item[Agent] A language model given tools and a loop, so it can act as well as
  write.

\item[Agent harness] The program that manages an agent's context, tools,
  permissions, memory and error recovery. These are where systems differ most.

\item[AGENTS.md / CLAUDE.md] A plain-text file of rules for agents working in a
  project, written in response to observed failures.

\item[AI-ready] Machine-readable, standardised, well-described, openly licensed
  data: good data management with a new consumer (\Cref{ch:aireadydata}).

\item[Attention] An operation that combines information from permitted
  positions in a sequence using calculated weights. Next-token prediction
  blocks access to later tokens \parencite{vaswani2017attention}.

\item[Balanced accuracy] The mean recall across classes, giving common and
  rare classes equal weight (\Cref{ch:applications}).

\item[Baseline] A simple method, such as keyword rules, run on the same items
  as the model, so that the model's score has something to be compared with
  (\Cref{ch:testing}).

\item[Benchmark contamination] A test item that was in a model's training
  data, so that a score measures memory rather than skill. Public reference
  databases and well-known papers are the usual route in. Invented items,
  which cannot have been memorised, are one control (\Cref{ch:testing}).

\item[CHOP (Chat-Oriented Programming)] Writing code in dialogue with a model,
  checking each step. Unlike vibe coding, checking is in the loop.

\item[Chunking] Splitting documents into passages for retrieval.

\item[Confidence interval] The range of scores consistent with the items
  tested. With ten items the range is wide, so report it beside the score
  (\Cref{ch:testing}).

\item[Context window] Everything a model can see in one call, measured in
  tokens. Material in the middle of long contexts is used less reliably.

\item[Coordinate reference system (CRS)] The datum and projection that give
  coordinates their meaning. Record it with every set of coordinates; the same
  numbers mean different places under different systems.

\item[COPE] Comparable, Organized, Predictive, Engaged: the work of harmonising
  and updating FAIR data so that it can be modelled
  \parencite{woodcharlson2026cope}.

\item[Coverage] The proportion of test cases for which a system supplies an
  answer. Report it alongside accuracy when a system can abstain.

\item[Cyberbiosecurity] The vulnerabilities that arise where digital systems
  and biological ones meet, from laboratories to public biological databases
  and DNA synthesis \parencite{ranjan2025cyberbiosecurity}
  (\Cref{ch:security}).

\item[Darwin Core] The standard vocabulary for sharing biodiversity records,
  maintained by Biodiversity Information Standards (TDWG). Its terms carry the
  fields, such as occurrence status and sampling effort, that make a record
  fit for analysis \parencite{wieczorek2012darwincore} (\Cref{sec:fitformodel}).

\item[Decode] See \emph{Prefill and decode}.

\item[Deep research] A search process that breaks a question into parts,
  reads the retrieved material and decides what to search next before
  bringing the results together.

\item[Dense model] A model in which every parameter is used for every token.
  Compare \emph{Mixture-of-experts}.

\item[Deskilling] Losing a skill by handing it to a tool, and with it the
  ability to check the tool \parencite{mammides2026llmecology}.

\item[Digital sequence information (DSI)] Sequence data derived from genetic
  resources, mostly held in open databases. Since 2024 its commercial use
  falls under a multilateral benefit-sharing mechanism, the Cali Fund
  \parencite{cbd2024dsidecision} (\Cref{sec:abs}).

\item[Embedding] A vector representing a text, such that similar texts lie close
  together. The basis of vector search. A text embedding model is a small,
  separate model that produces them; it is not the chat model.

\item[ENVO] The Environment Ontology: controlled terms for describing
  environments, used in habitat and sample fields (\Cref{ch:aireadydata}).

\item[FAIR] Findable, Accessible, Interoperable, Reusable: data principles
  biodiversity informatics followed long before ``AI-ready''.

\item[False zero] A zero or an absence written where nothing was counted: a
  lost sample, an unsearched site or an empty field read as a count. It enters
  an analysis as a true absence (\Cref{ch:testing}).

\item[F1 score] The harmonic mean of precision and recall:
  twice their product divided by their sum. Report the averaging method when
  combining classes or species.

\item[Fine-tuning] Further training of a model on domain data. Unlike RAG, it
  changes the weights.

\item[GBIF] The Global Biodiversity Information Facility: the aggregator of
  occurrence records published in Darwin Core, with the taxonomic backbone
  they are matched against.

\item[Gold set] A small set of reference items with known answers, used to
  measure what a search or a pipeline recovers (\Cref{ch:weeklyliterature}).

\item[Held-out set] Test items set aside before any rule, prompt or baseline
  is tuned, and used only to measure. The items used while tuning are the
  development set (\Cref{ch:testing}).

\item[Hallucination] Fluent, specific and false output. It persists because
  training rewards guessing over abstaining \parencite{kalai2025hallucinate}.

\item[I-ADOPT] A framework that breaks a measured variable into its parts,
  such as the property, the object and the unit, so that variables from
  different sources can be compared (\Cref{ch:aireadydata}).

\item[Inference] Two meanings. In statistics, drawing conclusions about a
  population from a sample. In machine learning, running a trained model to
  produce an output, including when you run a model locally. The book uses
  both: \Cref{ch:reproducibility} means the first in ``chain of inference'',
  and \Cref{ch:hardware} the second.

\item[Instruction tuning] The second training stage: fitting a pretrained
  model to example requests and answers, so that it follows instructions and
  declines some requests. See \emph{Training stages}.

\item[Jagged intelligence] Capability that is uneven and unpredictable, strong on
  hard tasks and failing on easy ones \parencite{mitchell2026jagged}.

\item[Jailbreaking] A user talking a model out of its safety training, so that
  it does what it was trained to refuse. Not the same as prompt injection,
  where content the model reads hijacks an application (\Cref{ch:security}).

\item[KV cache] Memory holding the keys and values computed for every token in
  the context, so that each new token does not recompute them. Its size per
  token is twice the product of the number of layers, the key-value heads,
  the head dimension and the bytes per value. It grows linearly with context
  and takes a large share of memory at long contexts (\Cref{ch:hardware}).

\item[Lethal trifecta] An agent that has private data, reads untrusted content
  and can send data out can be tricked into leaking the data. The safe course
  is never to combine all three \parencite{willison2025trifecta}
  (\Cref{ch:security}).

\item[Least privilege] Giving an agent only the permissions its current task
  needs, with explicit approval for anything more (\Cref{ch:security}).

\item[LLM] Large language model: a function from a token sequence to
  probabilities for the next token.

\item[Local-first] Software whose data lives on the user's device and works
  offline, with syncing optional \parencite{kleppmann2019localfirst}.

\item[Long-tail data] The many small, varied datasets held by individuals and
  small teams: together a substantial share of scientific data, and the part
  most often mismanaged or overlooked for want of curation
  \parencite{stahlman2026longtail} (\Cref{ch:organisation}).

\item[LoRA (Low-Rank Adaptation)] Fine-tuning a small set of added weights
  rather than the whole model. Cheap enough for one GPU.

\item[MaaS (Model as a Service)] Vendor-managed deployment of open-weight
  models, between hosted services and self-hosting.

\item[MCP (Model Context Protocol)] A standard way to offer tools and data to
  models. Tool descriptions are untrusted input.

\item[Memorisation] See \emph{Benchmark contamination}.

\item[MIxS] Minimum Information about any (x) Sequence: the Genomic Standards
  Consortium's checklists for describing the sample and context a sequence
  came from (\Cref{ch:aireadydata}).

\item[Mixture-of-experts (MoE)] A model whose layers hold several blocks of
  parameters, of which a router uses a few for each token. The total
  parameters set the memory a model needs; the active parameters set the
  compute per token, and so the speed. A tag such as 30B-A3B gives both
  (\Cref{ch:hardware}).

\item[Multi-token prediction] A feature built into some models: a cheaper
  mechanism inside the model guesses several tokens ahead. Runtimes use the
  guesses for speculative decoding, in which the full model checks them and
  keeps those that hold \parencite{willison2026qwen38}.

\item[Nagoya Protocol] The protocol to the Convention on Biological Diversity
  that sets out the procedures, permits and checkpoints for access and
  benefit-sharing \parencite{scholz2022dsi}.

\item[Non-detection] A record that a species was searched for and not found,
  with the effort spent. Without it, absence cannot be told from no search.

\item[OBIS] The Ocean Biodiversity Information System: the marine counterpart
  of GBIF.

\item[Occurrence record] One observation or specimen of a taxon at a place and
  time, the unit GBIF and OBIS share. Fit for analysis only with its method,
  effort and uncertainty.

\item[Open weight] Downloadable parameters. Not open source, which would also
  need the training data and code \parencite{hai2026openweight}.

\item[Orchestrator] The service between interface and model holding an
  application's rules: query expansion, retrieval, filters, citation checks.

\item[Parameter] One of the numbers a model learned in training. ``8B'' means
  eight billion of them; at 16 bits each, that is 16 GB of memory before
  quantisation.

\item[Phenomenal consciousness] Subjective experience, distinct from task
  performance or access to information for reasoning
  \parencite{chandaria2026consciousness}.

\item[Plazi treatment] The part of a taxonomic publication that describes one
  taxon, extracted and made citable on its own by Plazi
  (\Cref{ch:localfirst}).

\item[Precision] The fraction of positive predictions that are correct.
  In extraction, the fraction of extracted items matching the reference.

\item[Preference tuning] The third training stage: adjusting a model towards
  the answers people rate highly, often by reinforcement learning. The usual
  explanation of sycophancy. See \emph{Training stages}.

\item[Prefill and decode] The two phases of a model call. Prefill, or prompt
  processing, handles the whole prompt at once and is limited by compute;
  decode, or token generation, produces the answer one token at a time and
  is limited by memory bandwidth. Agent sessions with long prompts spend
  much of their time in prefill (\Cref{ch:hardware}).

\item[Pretraining] The first training stage: fitting a model to predict the
  next token over a large body of text.

\item[Prompt caching] Reusing the work done on a prompt prefix that repeats
  across calls, such as the rules file and the earlier turns of a session,
  instead of paying for it again.

\item[Prompt injection] Instructions hidden in content a model reads, which it
  then follows. Unsolved, because instructions and data share one channel.

\item[Provenance] The documented origin and processing history of a result.
  Missing from generated output unless you log it.

\item[Quantisation] Storing weights at lower precision to save memory; Q4 uses
  about a quarter of the memory of FP16. A tag such as \code{Q4\_K\_M} names
  the scheme: four bits on average, with a suffix for the variant, and
  variants differ in which layers keep more precision.

\item[RAG (Retrieval-Augmented Generation)] Supplying retrieved documents in the
  prompt. Reduces but does not remove hallucination
  \parencite{kurniawan2024ragfailures}.

\item[Reasoning effort] How long a model may reason privately before answering.
  Start low.

\item[Reasoning model] A model trained to write out intermediate steps before
  its answer, usually by reinforcement learning on tasks with checkable
  results. It spends more tokens per answer; the reasoning effort setting
  limits how many.

\item[Recall] The fraction of reference positives recovered by the system.
  For a literature search, the fraction of relevant reference papers retrieved.

\item[Reference-database version] The release of a taxonomic backbone or
  sequence database used in a step, such as the date of a GBIF backbone.
  Matches and scores change between releases, so record it with the result
  (\Cref{ch:testing}).

\item[Reinforcement learning on verifiable tasks] Training in which a model is
  rewarded when its answer can be checked by a program or a test, such as
  code that passes its tests. The method behind reasoning models
  (\Cref{ch:landscape}).

\item[Reranking] Re-scoring retrieved passages with a more careful model. Cheap
  and effective.

\item[RO-Crate] A standard bundle of research files and the metadata that
  describes them, including how they were produced.

\item[Sampling effort] How much searching produced a record: traps, hours,
  area, visits. Counts and absences mean little without it.

\item[Sandbox] An enclosure that the operating system or the harness puts
  around a process, limiting which files and network addresses it can reach.
  The enforcement layer that does not depend on the model following rules
  \parencite{anthropic2025sandboxing} (\Cref{ch:security}).

\item[Scaling laws] Regularities by which a model's error falls as a power law
  as its size, training data and compute grow. How best to divide a budget
  between size and data is disputed
  \parencite{kaplan2020scaling,hoffmann2022chinchilla} (\Cref{ch:whynow}).

\item[Seed] A number that fixes a random choice so that it can be repeated.
  It makes repeated runs closer, not identical (\Cref{ch:reproducibility}).

\item[Skill] A folder of instructions, and sometimes scripts, that an agent
  harness loads when a task matches its description or when a person invokes
  it by name. A procedure written once and reused across sessions
  (\Cref{ch:weeklyliterature}).

\item[Slop] Text produced in bulk by a model and passed on with little human
  thought in it \parencite{turner2026slop}.

\item[Slopsquatting] Registering package names that models invent, so generated
  install commands fetch an attacker's code.

\item[Speculative decoding] A faster way to generate: something cheaper than
  the full model (a small draft model, or the model's own multi-token
  prediction) proposes several tokens, and the full model checks them in one
  pass. The output is the full model's.

\item[Sycophancy] Excessive agreement with the user, usually linked to
  preference tuning, which rewards answers people rate highly
  \parencite{sharma2024sycophancy}. People trust a sycophantic model more and
  are more willing to return to it \parencite{cheng2026sycophancy}, and it can
  lead even rational users to false beliefs \parencite{chandra2026spiraling}.

\item[Taxonomic backbone] The reference list of names a database matches
  records against, such as the GBIF backbone, Catalogue of Life or WoRMS. It
  changes with each release; record which one and when.

\item[Temperature] A setting that controls how far a model's choice of the
  next token departs from the most probable one. Zero takes the most probable
  token each time (\Cref{ch:reproducibility}).

\item[Time-sharing] One computer divided among many users at terminals; the
  1960s answer to costly machines (\Cref{ch:landscape}).

\item[Token] A chunk of text, about four characters or three-quarters of a
  word in English prose. The unit of context and of billing. Scientific
  names, code and text in other languages split into more tokens per word.

\item[Training stages] Pretraining on a large body of text; instruction tuning
  on example requests and answers; preference tuning towards the answers
  people rate highly. Reasoning models add reinforcement learning on
  verifiable tasks \parencite{ouyang2022instructgpt,guo2025deepseekr1}
  (\Cref{ch:landscape}).

\item[Transformer] The architecture behind current language models (2017),
  using attention to process many positions in parallel during training.
  Its original decoder generates output one token at a time
  \parencite{vaswani2017attention}.

\item[Unit of replication] The thing that varies independently in a test,
  such as one field note; the cells of one note are not independent of each
  other. Intervals are computed over these units (\Cref{ch:testing}).

\item[Unified memory] Processor and graphics sharing one memory pool, so large
  models fit on ordinary machines.

\item[Validation sample] A random or stratified set of records that people
  check against the sources, to measure how often a model's output is wrong
  in each field (\Cref{sec:outputasdata}).

\item[Vibe coding] Accepting generated code without reading it. Fine for
  throwaway work only.

\end{description}
